\section*{第 \ref{ch:g11:color-light-sources} 章——颜色与光源}

\begin{solution}{exo:g11:color-light-sources:1}
（a）黄色；（b）青色；（c）白色。与蓝色互补的颜色：黄色
（蓝色 $+$ 黄色 $=$ 蓝色 $+$ 红色 $+$ 绿色 $=$ 白色）。
\end{solution}

\begin{solution}{exo:g11:color-light-sources:2}
黄色：R、G。品红：R、B。白色：R、G、B。黑色：无。橙色：R 取满\omterm{def:g10:energy-conservation:power}{功率}，
G 取降低的\omterm{def:g10:energy-conservation:power}{功率}。
\end{solution}

\begin{solution}{exo:g11:color-light-sources:3}
品红吸收绿色波段：$(R, B)$ 射出，为品红光。
青色\omterm{def:g11:color-light-sources:subtractive}{滤光片}随后吸收红色波段，只剩 $(B)$——蓝色。
\end{solution}

\begin{solution}{exo:g11:color-light-sources:4}
黄色条纹：\omterm{def:g11:color-light-sources:objectcolor}{漫射} R 与 G，吸收 B。（a）红光下它\omterm{def:g11:color-light-sources:objectcolor}{漫射} R，
呈红色；（b）蓝光下全部吸收，呈黑色。
\end{solution}

\begin{solution}{exo:g11:color-light-sources:5}
$T = \dfrac{\qty{2.90e-3}{m.K}}{\qty{5.80e-7}{m}} = \qty{5000}{K}$.
\end{solution}

\begin{solution}{exo:g11:color-light-sources:6}
\emph{(a)} $\lambda_{\max} = \qty{2.90e-3}{m.K}/\qty{1800}{K}
= \qty{1.61e-6}{m} = \qty{1610}{nm}$：\omterm{def:g10:light-spectra:wavelength}{红外线}。

\emph{(b)} \omterm{def:g10:light-spectra:white}{光谱}为连续谱：其可见部分很小，但并非
零，蓝端先减弱——留下的可见光主要是
红橙色，这便是火焰的颜色。
\end{solution}

\begin{solution}{exo:g11:color-light-sources:7}
外套（\omterm{def:g11:color-light-sources:objectcolor}{漫射} R）：(a) 红色；(b) 红色；(c) 黑色。手套（\omterm{def:g11:color-light-sources:objectcolor}{漫射} G 和
B，吸收 R）：(a) 青色；(b) 黑色；(c) 绿色。
\end{solution}

\begin{solution}{exo:g11:color-light-sources:8}
\emph{(a)} 为\omterm{def:g10:light-spectra:emission}{发射光谱}。每个物体只能\omterm{def:g11:color-light-sources:objectcolor}{漫射}
$\qty{589}{nm}$ 的光或不\omterm{def:g11:color-light-sources:objectcolor}{漫射}任何光：街道呈黄橙色的各种深浅，以及
黑色。

\emph{(b)} 蓝色汽车只\omterm{def:g11:color-light-sources:objectcolor}{漫射}蓝色波段。日光含有
它：蓝色汽车。钠谱线不含蓝光：全部被吸收，
汽车看上去呈黑色。
\end{solution}

\begin{solution}{exo:g11:color-light-sources:9}
\emph{(a)} \omterm{def:g10:light-spectra:white}{白光}束散开成彩虹；激光束
发生偏折，但仍是同色绿光的单束光：它是单色的
——只有一个\omterm{def:g10:light-spectra:wavelength}{波长}，棱镜无法将其分开。

\emph{(b)} 红色\omterm{def:g11:color-light-sources:subtractive}{滤光片}只透过红色波段，并吸收
$\qty{532}{nm}$：几乎没有光射出，光斑消失。
\end{solution}

\begin{solution}{exo:g11:color-light-sources:10}
\omterm{def:g10:light-spectra:white}{光谱}：在 $\qty{450}{nm}$ 附近有一个窄的蓝色峰，加上\omterm{def:g11:color-light-sources:led}{荧光粉}的
宽黄色波段（约 $500$--$\qty{700}{nm}$）。蓝色与黄色
互补，因此三种\omterm{def:g11:color-light-sources:cones}{视锥细胞}的响应\omterm{def:g11:color-light-sources:perception}{感知}为白色。该波段向
深红端减弱，因此深红色物体能\omterm{def:g11:color-light-sources:objectcolor}{漫射}的光很少：它看起来
比在日光下暗淡，日光的\omterm{def:g10:light-spectra:continuous}{连续光谱}能充分供给它。
\end{solution}

\begin{solution}{exo:g11:color-light-sources:11}
(a) $\lambda_{\max} = \qty{2.90e-3}{m.K}/\qty{1800}{K}
= \qty{1610}{nm}$，\omterm{def:g10:light-spectra:wavelength}{红外线}；
(b) $\qty{2.90e-3}{m.K}/\qty{6500}{K} = \qty{446}{nm}$，蓝紫色。
烛光被称为“暖色”——即温度
\emph{较低}的光源：称谓习惯与\omterm{def:g10:light-spectra:kelvin}{开尔文}温标相反。
\end{solution}

\begin{solution}{exo:g11:color-light-sources:12}
\emph{(a)} \omterm{def:g11:color-light-sources:perception}{同色异谱}。其存在证明颜色是一种\omterm{def:g11:color-light-sources:perception}{感知}
——三种\omterm{def:g11:color-light-sources:cones}{视锥细胞}的响应——而非\omterm{def:g10:light-spectra:white}{光谱}本身：无穷多
种\omterm{def:g10:light-spectra:white}{光谱}对应于同一组三元响应。

\emph{(b)} 棱镜或\omterm{def:g10:light-spectra:white}{分光镜}：在 $\qty{589}{nm}$ 处有一条谱线，对比
与 $540$ 和 $\qty{610}{nm}$ 处的两条波段相比。红色\omterm{def:g11:color-light-sources:subtractive}{滤光片}：它吸收
$\qty{589}{nm}$，钠光斑变暗，而屏幕的
$\qty{610}{nm}$ 子像素保留——屏幕上的该色块变为红色。
\end{solution}

\begin{solution}{exo:g11:color-light-sources:13}
\emph{(a)} $\lambda_{\max} = \qty{2.90e-3}{m.K}/\qty{3200}{K} =
\qty{906}{nm}$。\omterm{def:g11:color-light-sources:colortemp}{色温} $\qty{3200}{K} < \qty{6500}{K}$：
比日光更暖（更偏橙色）。

\emph{(b)} 滤色片吸收多余的部分红橙光，以使
\omterm{def:g10:light-spectra:white}{光谱}向蓝端重新平衡；由于\omterm{def:g11:color-light-sources:subtractive}{滤光片}只会去除光，舞台上的总
\omterm{def:g10:energy-conservation:power}{功率}下降。

\emph{(c)} \omterm{def:g11:color-light-sources:subtractive}{滤光片}是被动的：它可以在每个\omterm{def:g10:light-spectra:wavelength}{波长}处吸收，但
从不发光，因此没有任何色片能使蓝光光强高于光斑
本身所产生的强度。
\end{solution}

\begin{solution}{exo:g11:color-light-sources:14}
\emph{(a)} $\qty{2700}{K}$：暖色、富红光——肉的红色波段
得到充分供给，肉看起来鲜艳新鲜。$\qty{6500}{K}$：冷色
偏蓝的\omterm{def:g10:light-spectra:white}{白光}——鱼完整地\omterm{def:g11:color-light-sources:objectcolor}{漫射}三色光，看起来明亮
洁白，新鲜如刚冰镇过。

\emph{(b)} 两者的货品都未改变：肉与鱼的吸收特性
并未改变。改变的是照明，因而\omterm{def:g11:color-light-sources:objectcolor}{漫射}
的光——\omterm{def:g11:color-light-sources:objectcolor}{物体的颜色}只在给定
照明条件下才有确切的意义。
\end{solution}

\begin{solution}{exo:g11:color-light-sources:15}
\emph{(a)} 钠灯：全城的光晕集中于一条谱线，位于
$\qty{589}{nm}$ 处，一块窄带\omterm{def:g11:color-light-sources:subtractive}{滤光片}即可将其从望远镜
图像中滤除；宽带的 LED \omterm{def:g10:light-spectra:white}{光谱}则无法滤除，
否则星光也会一并损失。

\emph{(b)} 司机和行人更喜欢 LED：在其完整\omterm{def:g10:light-spectra:white}{光谱}下，
\omterm{def:g11:color-light-sources:objectcolor}{物体的颜色}仍可分辨（红色外套依旧是红的），
而钠光则使一切都呈橙黄色或黑色。

\emph{(c)} 暖白色 LED，\omterm{def:g11:color-light-sources:colortemp}{色温}约 $\qty{2700}{K}$：显色性尚可，
适合街道照明；其\omterm{def:g10:light-spectra:white}{光谱}蓝端较弱，而蓝光
恰是对夜空污染最严重的部分。
\end{solution}

\begin{solution}{pb:g11:color-light-sources:1}
\textbf{1.} R $+$ G $=$ 黄色，R $+$ B $=$ 品红，G $+$ B $=$ 青色；
三色重叠为白色。

\textbf{2.} 橙色：红光满\omterm{def:g10:energy-conservation:power}{功率}，绿光部分\omterm{def:g10:energy-conservation:power}{功率}，蓝光
关闭。浅粉色：三色均开（以白色为基础），红光略占优势。

\textbf{3.} 红色光斑投下的影子只接收到绿光光斑的光，
因而为绿色影子；对称地，另一个影子为红色。

\textbf{4.} 深蓝色涂料只\omterm{def:g11:color-light-sources:objectcolor}{漫射}蓝光。仅在红光光斑下：
黑色。三色光斑同时照射时：它\omterm{def:g11:color-light-sources:objectcolor}{漫射}蓝光波段——呈蓝色。

\textbf{5.} 眼睛将任何\omterm{def:g10:light-spectra:white}{光谱}都简化为三种\omterm{def:g11:color-light-sources:cones}{视锥细胞}的响应；
只要三者匹配，颜色便相同。三种\omterm{def:g11:color-light-sources:cones}{视锥细胞}决定了
所需聚光灯的数目。

\textbf{6.} 黄色 $=$ \omterm{def:g10:light-spectra:white}{白光} $-$ 蓝光：染料吸收蓝光波段。

\textbf{7.} 红色光斑：\omterm{def:g11:color-light-sources:objectcolor}{漫射} R——呈红色。绿色光斑：\omterm{def:g11:color-light-sources:objectcolor}{漫射} G——
呈绿色。蓝色光斑：吸收一切——呈黑色。

\textbf{8.} 仅用蓝色光斑照亮舞台；灯光由
白色 $\to$ 蓝色。黄色戏服\omterm{def:g11:color-light-sources:objectcolor}{漫射} R 和 G，却吸收 B：
\omterm{def:g10:light-spectra:white}{白光}下明亮，蓝光下呈黑色。

\textbf{9.} 品红色\omterm{def:g11:color-light-sources:subtractive}{滤光片}：$(R, G, B) \to (R, B)$。戏服\omterm{def:g11:color-light-sources:objectcolor}{漫射}
所接收光线中的 R 与 G：它\omterm{def:g11:color-light-sources:objectcolor}{漫射} R、吸收 B——红色戏服。

\textbf{10.} 红色\omterm{def:g11:color-light-sources:subtractive}{滤光片}只保留三个波段中的一个，
白色光斑的\omterm{def:g10:energy-conservation:power}{功率}因此损失三分之二；每块青、品红、黄\omterm{def:g11:color-light-sources:subtractive}{滤光片}
各去除一个波段，保留三分之二，两块叠加仍可得到任一原色。

\textbf{11.} \omterm{def:g10:light-spectra:emission}{发射光谱}。街道变成单色的世界：
每个表面都呈橙黄色，只是明暗不同，或呈黑色。

\textbf{12.} 白墙：橙黄色（\omterm{def:g11:color-light-sources:objectcolor}{漫射}该谱线）。黄色
标线：明亮的橙黄色。红色汽车：暗色——$\qty{589}{nm}$
谱线位于其所\omterm{def:g11:color-light-sources:objectcolor}{漫射}的波段之外。蓝色汽车：黑色。

\textbf{13.} 中心位于 $\qty{589}{nm}$ 的一块窄带\omterm{def:g11:color-light-sources:subtractive}{滤光片}，可从图像中去除
整个城市的光晕，同时几乎不损失
恒星的宽带\omterm{def:g10:light-spectra:white}{光谱}。

\textbf{14.} $\lambda_{\max} = \qty{2.90e-3}{m.K}/\qty{2700}{K} =
\qty{1.07e-6}{m} = \qty{1070}{nm}$：峰值及大部分\omterm{def:g10:energy-conservation:power}{功率}
位于\omterm{def:g10:light-spectra:wavelength}{红外线}波段——是热量，而非光。

\textbf{15.} 钠灯：$0.30 \times \qty{100}{W} = \qty{30}{W}$ 的
可见光；钨丝灯：$0.05 \times \qty{100}{W} = \qty{5}{W}$。六
倍的光，电费相同。

\textbf{16.} \omterm{def:g11:color-light-sources:led}{荧光粉}将部分蓝光转换为宽带黄光，
波段；蓝、黄互为\omterm{prop:g11:color-light-sources:complementary}{互补色}，因此蓝光 $+$ 黄光
即合成\omterm{def:g10:light-spectra:white}{白光}。

\textbf{17.} 一个狭窄的蓝色尖峰加上一个宽广的黄色驼峰——与
日光均匀连续的\omterm{def:g10:light-spectra:white}{光谱}相比。深红端（以及蓝、黄之间的低谷）
明显偏弱。

\textbf{18.} 深红色物体显色暗淡，偏棕色。因此肉铺
采用 $\qty{2700}{K}$ 暖色 LED，其增强的红色成分
使肉色保持鲜艳。

\textbf{19.} 内部并非处于 $\qty{2700}{K}$：芯片工作时接近室温。
标签只承诺\emph{色调}与一个
$\qty{2700}{K}$ \omterm{def:g11:color-light-sources:incandescent}{白炽}体相同，其峰值将位于
$\qty{2.90e-3}{m.K}/\qty{2700}{K} = \qty{1.07e-6}{m} = \qty{1070}{nm}$。

\textbf{20.} 卧室：约 $\qty{2700}{K}$——暖色、少蓝光的光线
利于休息。珠宝店：$5500$--$\qty{6500}{K}$——类似日光的\omterm{def:g10:light-spectra:white}{白光}，使
白色看起来仍为白色，宝石能反射每一个波段。规律：所\omterm{def:g11:color-light-sources:perception}{感知}的
颜色 $=$ 光源\omterm{def:g10:light-spectra:white}{光谱} $\times$ 物体吸收，经
三种\omterm{def:g11:color-light-sources:cones}{视锥细胞}读出。
\end{solution}
